\section{UML модели}

\begingroup

\textit{Скоро тут будет концептуальная диаграмма классов ВКР.
Пока диаграммы добавлены сюда для примера рисования UML средствами Tikz. Приложение недопридумано, недоделано и недооформлено. }


%Все UML-либы не очень. 
%Везде всё не верно, я разрабам даже 25 бы не поставил по СРСППО (или ТРПО). Придется делать все самому.

\def\baselinestretch{1}

\def\umlDebug{}
	
%#1 Tikz Опции
%#2 стереотип
%#3 имя нода 
%#4 имя класса
%#5 атрибуты
%#6 методы
\NewDocumentCommand{\class}{O{}d<>mD(){clstmp}mm}{
	\node[class,
	#1] 
	(#4)
	{%   
		\centering%
		\IfNoValueF{#2}{\textit{\phantom{рб}#2}\\}%
		\textbf{\strut#3}%
		\par	%		
		%
		\IfBlankF{#5}{%
			\nodepart{two}%
			#5%
		}%
		\IfBlankF{#6}{%
			\nodepart{three}%
			#6%
		}%
	};
	\ifdefined\umlDebug%
	\draw[red] (#4.north east) node[font=\fontsize{8pt}{5pt}\selectfont, inner sep=0pt, outer sep=1pt, above left] {#4};%
	\fi
}

%рисует линию со ступенькой
%создает координаты a и b  на начало и конец
%1 - свойства линии Tikz
%2 - начало
%3 - конец
%4 - горизонтальное расстояние от начала до ступеньки
\NewDocumentCommand{\steplineH}{O{}mmm}{
	\draw [#1] (#2)coordinate (a) -- ++(#4,0)  |- (#3) coordinate(b);
}

%Аналог для вертикальной линии
\NewDocumentCommand{\steplineV}{O{}mmm}{
	\draw [#1] (#2)coordinate (a) -- ++(0,#4)  -| (#3) coordinate(b);
}

\tikzset{
	own/.tip={Circle[]},
	base/.tip={Triangle[open,scale width=1,scale length=1,scale=2]},
	agr/.tip={Diamond[open,scale=2.5]},
	comp/.tip={Diamond[scale=2.5]},
	nav/.tip={Straight Barb[scale length=1, scale width=1,scale=1.6]_[sep=-0.1pt]},
	oagr/.tip={.agr[]own[]},
	ocomp/.tip={comp[]own[]},
	onav/.tip={nav[]own[]},
	ocompnav/.tip={nav[]ocomp[]},
	oagrnav/.tip={nav[]oagr},
	compnav/.tip={nav[]comp[]},
	agrnav/.tip={nav[].agr[]},
	lnk/.style={
		align=center, 
		font=\fontsize{8pt}{10pt}\selectfont,
		inner sep=1.3ex,
	},
	class/.style={
		font=\fontsize{10pt}{12pt}\selectfont,
		rectangle split, 
		rectangle split parts=3, 
		rectangle split ignore empty parts,
		draw,
		text width=3cm,
		rectangle split part align={left, left, left},
		inner sep=2pt,
	},
}






\centering


\begin{tikzpicture}
	
	%Класс документ
	\class[]{\itshape Документ}(cDoc){}{}
	
	%класс Элемент
	\class[right = 4cm of  cDoc.north east, text depth = 3cm, below right]{\itshape Элемент}(cEl){}{}
	
	%Наследование Элемент - Документ
	\draw[base-] (cEl.west) ++(0,-0.5) -| (cDoc.south);
	
	%Класс Комплект документов
	\class[above = of cDoc, text width = 4cm]{\itshape КомплектДокументов}(cDocPack){}{}
	
	%Агрегация Комплект-документ
	\steplineH[agr-onav]{cDocPack.west}{cDoc.west}{-1.5cm}
	\draw (a)
	node[lnk, above left, align=right]{+комплект}
	node[lnk, below left, align=right]{*};
	\draw (b)
	node[lnk, above left, align=right]{+документ}
	node[lnk, below left, align=right]{*};
	
	
	
	%Наследовние Элемент - Комплект
	\draw[base-] (cEl.north west) ++ (1,0) |- (cDocPack);
	
	
	%Композиция Документ-Элемент
	\draw (cDoc.east) coordinate (a)  [compnav-onav]-- (a -| cEl.west) coordinate (b);	
	\draw (a) 
	node[lnk, above right, align=left] {\{union, readOnly,\\subsets базовый\}\\+/документ}
	node[lnk, below right] {\quad 1}; 
	\draw (b)
	node[lnk, below left, align=right] {\{union, readOnly,\\subsets вложенный\}\\+/элемент} 
	node[lnk, above left] {0..*\quad}; 

	
	%Композиция Элемент-элемент
	\steplineH[ocompnav-onav]{$(cEl.east)+(0,1)$}{$(cEl.east)-(0,1)$}{2.5cm}
	\draw (a) 
	node[lnk, above right,align=left] {\{union, readOnly\}\\+/базовый} 
	node[lnk, below right,align=left] {\quad 0..1};	
	\draw (b)
	node[lnk, above right,align=left] {\{union, readOnly\}\\+/вложенный} 
	node[lnk, below right,align=left] {\quad *};
	
	%класс Ссылка
	\class[below = 2cm of cDoc.south east, below left]
 	{\itshape Ссылка}(cLink){}{};
 	
 	%Наследование Элемент--Ссылка
 	\draw[-base] (cLink.east)  coordinate (a) -- (a -| cEl.west);
 	
 	%Связь Ссылка - Элемент
 	\steplineV[-onav]{cLink.south}{$(cEl.south)-(1,0)$}{-1.5cm};
 	\draw (a) node[lnk, below right, align=left] {\{union, readOnly\}\\+/источник}
 		      node[lnk, below left, align=left] {*}
 		  (b) node[lnk, below left, align=right] {\{union, readOnly\}\\+/цель}
 		  	  node[lnk, below right, align=right] {0..1};
			
\end{tikzpicture}


\vspace{1cm}\rule{\linewidth}{1pt}\vspace{1cm}


\begin{tikzpicture}


	%класс документ
	\class[]{\itshape Документ}(cDoc){}{}
	
	%класс ПростойДокумент
	\class[text width=5cm, below = of cDoc]{\itshape ПростойДокумент}(cSDoc){}{}
	
	%наследование Документ-ПростойДокумент
	\draw [base-] (cDoc) -- (cSDoc);
	
	%Класс СодержательныйЭлемент
	\class[right = 4cm of cSDoc.north east, below right, text width=4cm ] {\itshape СодержательныйЭлемент}(cContEl){}{}
	
	%Класс Элемент
	\class[above = of cContEl] {\itshape Элемент}(cEl){}{}
	
	%наследование Элемент-Сод.Элемент
	\draw [base-] (cEl) -- (cContEl);
	
	%Композиция ПростойДокумент-Содэлемент	
	\draw[comp-onav] (cSDoc) -- (cContEl)
	node [lnk, pos=0, align=left, above right] {\{unon, readOnly,\\subsets документ\}\\+/простойДокумент}
	node [lnk, pos=0, align=left, below right] {\quad 1}
	node [lnk, pos=1, align=right, below left] {\{unon, readOnly,\\subsets элемент\}\\+/содЭлемент}
	node [lnk, pos=1, align=right, above left] {*\quad};
	
	
	
\end{tikzpicture}

\vspace{1cm}\rule{\linewidth}{1pt}\vspace{1cm}

\begin{tikzpicture}
	
	%Класс СодержательныйЭлемент
%	\class[right = 4cm of cSDoc.north east, below right, text width=4cm ] {\itshape СодержательныйЭлемент}(cContEl){}{}

	\node[matrix, column sep=0.2cm]
	{
		\class {БлокТекста} (сText){}{} \\
		\class {БлокРисунка} (сPic){}{} \\
		\class {БлокКода} (сText){}{}   \\
		\class {БлокТаблицы} (сText){}{} \\
		\class {КомбоБлок} (сText){}{}   \\
	};
	
\end{tikzpicture}

\vspace{1cm}\rule{\linewidth}{1pt}\vspace{1cm}

\begin{tikzpicture}
	
	%класс СтруктурированныйДокумент
	\class[text width=5cm]{\itshape СтруктурированныйДокумент}(cSDoc){}{}
	
	%класс Документ
	\class[above = of cSDoc]{\itshape Документ}(cDoc){}{}
	
	%Наследование Документ-СтруктурированныйДоукмент
	\draw[-base] (cSDoc) -- (cDoc);
	
	%Класс Часть
	\class[right = 4cm of cSDoc.north east, below right, text depth=3cm]{\itshape Часть}(cPart){}{}
	
	%Класс Элемент
	\class[above = of cPart]{\itshape Элемент}(cEl){}{}
	
	%Наследование Элемент - Часть
	\draw [base-] (cEl) -- (cPart);
	
	%композиция Часть-Часть
	\steplineH[ocompnav-onav]{$(cPart.east)+(0,1)$}{$(cPart.east)-(0,1)$}{2}
	\draw (a)
	node[lnk,above right, align=left] {\{union, readOnly\}\\+/контекст}
	node[lnk,below right] {\quad 0..1}
	(b)
	node[lnk,above right, align=left] {\{union, readOnly\}\\+/подчасть}
	node[lnk,below right] {\quad *};
	
	%Композиция СтруктурированныйДокумент - Часть
	\draw[ocompnav-onav](cSDoc.east) coordinate (a) -- (a -| cPart.west)
	node [lnk, pos=0, above right, align=left] {\{union, readOnly,\\subsets документ\}\\+/стрДокумент}
	node [lnk, pos=0, below right, align=left] {\quad 1}
	node [lnk, pos=1, below left, align=right] {\{union, readOnly,\\subsets элемент\}\\+/часть}
	node [lnk, pos=1, above left, align=left] {1..*\quad}; 
	
	
	
\end{tikzpicture}









%=======================================================
%Архив, исходные варианты невошедших диаграмм.
%=======================================================
\if0


\begin{tikzpicture}
	
	\class[text width=6cm]
	{\itshape ТекстовыйДокумент\\}(clsTxt)
	{+наименование:строка\vspace*{3cm}}{}
	
	\class[%
	below = 3cm of clsTxt,
	text width=6cm
	]
	{\itshape Часть}(clsPart)
	{}{}
	
	\draw [{comp[]nav[]}-onav] (clsTxt) -- (clsPart) 
	node[lnk, pos=0, align=left, below right] {\{union\}\\+документ}
	node[lnk, pos=1, above right,lnk,align=left] {\{union\}\\+/часть[1..*]};
	
	
	\class[right = 5cm of clsTxt.north east, below right]%
	{\textit{Автор}}(clsAuthor)
	{}{}	
	
	
	\draw[onav-onav] (clsAuthor) -- (clsAuthor.180 -| clsTxt.0)
	node[lnk, above left, pos=0] {\{union\}\\+/автор}
	node[lnk, below left, pos=0] {[1..*]}
	node[lnk, above right, pos=1] {\{union\}\\+/документ}
	node[lnk, below right, pos=1] {[0..*]};
	
	\class[below = 2cm of clsAuthor.south west, below right]%
	{\textit{Куратор}}(clsController)
	{}{}
	
	\draw[onav-onav] (clsController) -- (clsController.180 -| clsTxt.0)
	node[lnk, above left, pos=0] {\{union\}\\+/куратор}
	node[lnk, below left, pos=0] {[1..*]}
	node[lnk, above right, pos=1] {\{union\}\\+/документ}
	node[lnk, below right, pos=1] {[0..*]};
	
	\draw ($(clsPart.south)-(1,0)$) 
	node[lnk, below left, align=right] {\{union\}\\+/часть} [compnav-] -- 
	++(0,-1)  
	[-onav] -|  
	($(clsPart.south)+(1,0)$) 
	node[lnk, below right, align=left] {\{union\}\\+/подчасть[0..*]};
	
	\class[right = 2cm of clsPart.north east, text width = 5cm, below right]
	{\itshape ИменноваяЧасть}(clsNamedPart)
	{+наименование:строка}{}
	
	\draw [base-] (clsPart) -- ( clsPart.0 -| clsNamedPart.180);
	
\end{tikzpicture}

\newpage




\begin{tikzpicture}
	
	\class[text width=3cm]{\itshape Документ}(cDoc){}{}
	
	\class[right = 4cm of  cDoc.north east, text depth = 3cm, below right]{\itshape Элемент}(cEl){}{}
	
	\draw (cDoc.east) coordinate (a)  [compnav-onav]-- (a -| cEl.west) coordinate (b);	
	\draw (a) 
	node[lnk, above right] {\{union\}\\{+/документ}} 
	node[lnk, below right] {\quad[1]};	
	\draw (b)
	node[lnk, below left] {\{union\}\\{+/элемент}} 
	node[lnk, above left] {[0..*]\quad};
	
	\draw[ocompnav-onav] ($(cEl.east)+(0,1)$) coordinate (a)  -- +(2.5cm,0) |- ($(cEl.east)-(0,1)$) coordinate (b) ;	
	\draw (a) 
	node[lnk, above right,align=left] {\{union\}\\+/базовый} 
	node[lnk, below right,align=left] {\quad[1]};	
	\draw (b)
	node[lnk, above right,align=left] {\{union\}\\+/вложенный} 
	node[lnk, below right,align=left] {\quad[0..*]};
	
	\class[text width=6cm, below = 5cm of cDoc]{\itshape СтруктурированныйДокумент\strut}(cStrDoc){}{}
	
	\class[text width=3cm, text depth = 4.5cm, right = 5cm of cStrDoc.north east, below right]{\itshape \strut Часть}(cPart){}{}
	
	\draw ([xshift=1.2cm]cEl.south) coordinate (a) [base-]-- (a |- cPart.north);
	
	\draw (cDoc.south) coordinate (a) [base-]-- (a |- cStrDoc.north) coordinate (b);
	
	\draw (cStrDoc.east) coordinate (a) [compnav-onav]-- (a -| cPart.west) coordinate (b);
	\draw (a) 
	node[lnk, above right] {\{subsets документ\}\\{+/стр.документ}} 
	node[lnk, below right] {\quad[1]};	
	\draw (b)
	node[lnk, below left] {\{subsets элемент\}\\{+/часть}} 
	node[lnk, above left] {[0..*]\quad};
	
	\draw [onav-ocompnav] ($(cPart.south west)+(0,0.5cm)$) coordinate(a) -- ++(-4cm,0) |- ($(cPart.south west)+(0,2.5cm)$) coordinate(b);
	\draw (a) 
	node[lnk, above left,align=right] {\{union, subsets вложенный\}\\+/подчасть} 
	node[lnk, below left,align=right] {[0..*]\quad};	
	\draw (b)
	node[lnk, above left,align=right] {\{union, subsets базовый\}\\+/часть} 
	node[lnk, below left,align=right] {[1]\quad};
	
	%\class [right = 2cm of cEl.west]{Ссылка}
	
	
	
\end{tikzpicture} 

\fi


\endgroup
	
